\documentclass[12pt,a4paper]{article}

% ==================================================
% PACKAGES
% ==================================================

\usepackage[utf8]{inputenc}
\usepackage[T2A]{fontenc}
\usepackage[english,mongolian]{babel}

\usepackage{amsmath,amssymb,amsfonts}
\usepackage{mathtools}

\usepackage{graphicx}
\usepackage{booktabs}
\usepackage{array}
\usepackage{multirow}

\usepackage{geometry}
\geometry{
	left=2.5cm,
	right=2.5cm,
	top=2.5cm,
	bottom=2.5cm
}

\usepackage{setspace}
\onehalfspacing

\usepackage{natbib}

\usepackage{hyperref}
\hypersetup{
	colorlinks=true,
	linkcolor=black,
	citecolor=black,
	urlcolor=blue
}

% ==================================================
% TITLE
% ==================================================

\title{
	\textbf{A Comparison of Covariance Matrix Shrinkage Methods for Portfolio Optimization under a Low Asset-to-Observation Ratio}
}

\author{
	Narangoo Ganbaatar$^{1}$ \qquad Ankhbayar Chuluunbaatar$^{2}$\\
	\small $^{1}$School of Information Technology and Electronics, National University of Mongolia\\
	\small \texttt{narangoo.g@num.edu.mn}\\
	\small $^{2}$Business School, National University of Mongolia\\
	\small \texttt{ankhbayar.ch@num.edu.mn}
}

\date{}

\date{}

% ==================================================
% DOCUMENT
% ==================================================

\begin{document}
	
	\maketitle
	
	% ==================================================
	% ABSTRACT
	% ==================================================
	
	\begin{abstract}
		
		This study compares alternative covariance matrix shrinkage
		estimators in the context of portfolio selection.
		The performance of the estimators is evaluated using the
		Global Minimum Variance Portfolio (GMVP).
		
		\medskip
		
		\noindent
		\textbf{Keywords:}
		Covariance matrix,
		Shrinkage estimation,
		Portfolio selection,
		Global minimum variance portfolio,
		Ledoit--Wolf.
		
	\end{abstract}
	
	
	% ==================================================
	% 1. INTRODUCTION
	% ==================================================
	
	\section{Introduction}
	
	Portfolio optimization requires reliable estimates of the
	covariance matrix of asset returns.
	
	The traditional sample covariance matrix may perform poorly,
	particularly when the number of assets is large relative to
	the number of observations.
	
	Shrinkage estimators provide an alternative approach by
	combining the sample covariance matrix with a structured
	target matrix.
	
	The main objective of this study is to compare different
	covariance matrix shrinkage methods in the context of
	minimum-variance portfolio selection.
	
	
	% ==================================================
	% 2. LITERATURE REVIEW
	% ==================================================
	
	\section{Literature Review}
	
	Covariance matrix estimation has received considerable
	attention in portfolio theory.
	
	The shrinkage approach was introduced into portfolio
	applications by \citet{ledoit2004}.
	
	More recent studies have considered both linear and
	nonlinear shrinkage estimators.
	
	
	% ==================================================
	% 3. METHODOLOGY
	% ==================================================
	
	\section{Methodology}
	
	
	% --------------------------------------------------
	\subsection{Sample Covariance Matrix}
	% --------------------------------------------------
	
	Let
	
	\[
	R_t =
	\begin{pmatrix}
		R_{1t}\\
		R_{2t}\\
		\vdots\\
		R_{Nt}
	\end{pmatrix}
	\]
	
	denote the vector of asset returns.
	
	The sample covariance matrix is
	
	\begin{equation}
		S
		=
		\frac{1}{T}
		\sum_{t=1}^{T}
		(R_t-\bar R)(R_t-\bar R)'.
	\end{equation}
	
	
	% --------------------------------------------------
	\subsection{Linear Shrinkage Estimation}
	% --------------------------------------------------
	
	The general linear shrinkage estimator is defined as
	
	\begin{equation}
		\widehat{\Sigma}
		=
		\delta F
		+
		(1-\delta)S,
	\end{equation}
	
	where
	
	\begin{itemize}
		
		\item $S$ is the sample covariance matrix,
		
		\item $F$ is the shrinkage target matrix,
		
		\item $\delta \in [0,1]$ is the shrinkage intensity.
		
	\end{itemize}
	
	
	% --------------------------------------------------
	\subsection{Shrinkage Targets}
	% --------------------------------------------------
	
	Several target matrices are considered in this study.
	
	\subsubsection{Constant Variance Target}
	
	\subsubsection{Diagonal Target}
	
	\subsubsection{Constant Correlation Target}
	
	\subsubsection{Single-Factor Target}
	
	
	% --------------------------------------------------
	\subsection{Nonlinear Shrinkage}
	% --------------------------------------------------
	
	The nonlinear shrinkage estimator can be expressed as
	
	\begin{equation}
		\widehat{\Sigma}_{NLS}
		=
		U
		\widehat{D}
		U',
	\end{equation}
	
	where $U$ contains the eigenvectors of the sample covariance
	matrix and
	
	\[
	\widehat{D}
	=
	\operatorname{diag}
	(\widehat d_1,\ldots,\widehat d_N)
	\]
	
	contains the shrunk eigenvalues.
	
	
	% ==================================================
	% 4. PORTFOLIO OPTIMIZATION
	% ==================================================
	
	\section{Portfolio Optimization}
	
	This study considers the Global Minimum Variance Portfolio
	(GMVP).
	
	The optimization problem is
	
	\begin{equation}
		\min_{w}
		\quad
		w'\widehat{\Sigma}w
	\end{equation}
	
	subject to
	
	\begin{equation}
		\mathbf{1}'w = 1.
	\end{equation}
	
	The analytical solution is
	
	\begin{equation}
		w_{\mathrm{GMVP}}
		=
		\frac{
			\widehat{\Sigma}^{-1}\mathbf{1}
		}{
			\mathbf{1}'
			\widehat{\Sigma}^{-1}
			\mathbf{1}
		}.
	\end{equation}
	
	
	% ==================================================
	% 5. DATA
	% ==================================================
	
	\section{Data}
	
	This section describes the dataset used in the empirical
	analysis.
	
	\begin{table}[htbp]
		
		\centering
		
		\caption{Descriptive Statistics}
		
		\begin{tabular}{lrrrr}
			
			\toprule
			
			Asset & Mean & Std. Dev. & Minimum & Maximum \\
			
			\midrule
			
			Asset 1 & 0.00 & 0.00 & 0.00 & 0.00\\
			Asset 2 & 0.00 & 0.00 & 0.00 & 0.00\\
			Asset 3 & 0.00 & 0.00 & 0.00 & 0.00\\
			
			\bottomrule
			
		\end{tabular}
		
	\end{table}
	
	
	% ==================================================
	% 6. EMPIRICAL RESULTS
	% ==================================================
	
	\section{Empirical Results}
	
	
	\subsection{Covariance Matrix Estimation}
	
	The estimated covariance matrices are compared in terms of
	their statistical properties.
	
	
	\subsection{Portfolio Weights}
	
	The GMVP weights obtained from different covariance
	estimators are compared.
	
	
	\subsection{Portfolio Performance}
	
	Portfolio performance is evaluated using volatility,
	variance, Sharpe ratio, and portfolio turnover.
	
	\begin{table}[htbp]
		
		\centering
		
		\caption{GMVP Performance Comparison}
		
		\begin{tabular}{lrrrr}
			
			\toprule
			
			Estimator &
			Return &
			Volatility &
			Sharpe Ratio &
			Turnover \\
			
			\midrule
			
			Sample Covariance & & & & \\
			Constant Variance & & & & \\
			Diagonal & & & & \\
			Constant Correlation & & & & \\
			Single Factor & & & & \\
			Nonlinear Shrinkage & & & & \\
			
			\bottomrule
			
		\end{tabular}
		
	\end{table}
	
	
	% ==================================================
	% 7. DISCUSSION
	% ==================================================
	
	\section{Discussion}
	
	This section discusses the empirical differences between
	the covariance matrix estimators and their implications for
	portfolio selection.
	
	
	% ==================================================
	% 8. CONCLUSION
	% ==================================================
	
	\section{Conclusion}
	
	This study compares several covariance matrix shrinkage
	methods for constructing global minimum variance portfolios.
	
	The empirical results demonstrate how covariance estimation
	affects portfolio risk, asset allocation, and portfolio
	stability.
	
	
	% ==================================================
	% REFERENCES
	% ==================================================
	
	\nocite{*}
	\bibliographystyle{apalike}
	\bibliography{references}
	
\end{document}